\subsection{R-35: Behrend sets extend the Weyl-pencil range to nearly quadratic size}

\paragraph{Statement R-35.}
There are absolute constants $C>0$ and $m_0$ such that, for every
integer $m\geq m_0$, there is an integer $R_m$ satisfying
\[
 m^2\exp\!\bigl(-C\sqrt{\log m}\bigr)\leq R_m<m^2
\]
for which every conclusion of R-34 holds for every number of rows
$3\leq r\leq R_m$. In particular, for every fixed $\delta>0$, those
conclusions hold for all $3\leq r\leq m^{2-\delta}$ once $m$ is
sufficiently large. This includes arbitrary composite and even $m$.
The conclusions retain the scope of R-34: the permanent and determinant
witnesses may be different, and an unpadded rectangular direction is
asserted to have a stable point but is not asserted to have degree
period one.

\begin{proof}
We use the classical Behrend bound in the following equivalent
interval formulation: for some absolute constant $C_0>0$ and all
sufficiently large integers $L$, there is a set
$B\subset\{1,\ldots,L\}$ with no nontrivial three-term arithmetic
progression and
\[
                 |B|\geq L\exp\!\bigl(-C_0\sqrt{\log L}\bigr).
\]
This is the cited theorem of Behrend (1946, p.~331, second displayed
estimate, with its proof completed on p.~332), not a new density
estimate. Translating the interval and enlarging the constant gives
the displayed formulation from his convention $B\subset\{0,\ldots,N\}$.

Set
\[
 L=\left\lfloor\frac{m-5}{6}\right\rfloor,
 \qquad A=\{0,1\}\cup(2L+2+B).
\]
Then the translated part of $A$ lies in
$[2L+3,3L+2]$, and $A\subset[0,(m-1)/2]$. We check the stronger
modular condition
\[
 a+c=2b\pmod m,\quad a,b,c\in A
                 \quad\Longrightarrow\quad a=b=c.
\]
First, both integer quantities $a+c$ and $2b$ lie in $[0,m-1]$,
so the congruence is an integer equality. A progression lying
entirely in the translated part is trivial by the defining property
of $B$. A progression entirely in $\{0,1\}$ is also trivial.
For a mixed progression, if the middle term is in $\{0,1\}$, the
presence of a translated endpoint makes the endpoint sum exceed
$2$, which is impossible. If the middle term is in the translated
part, twice that term is at least $4L+6$, whereas an endpoint sum
with at least one endpoint in $\{0,1\}$ is at most $3L+3$.
This is again impossible. The argument includes equal endpoints;
there is no residual two-torsion exception for even $m$.

Consequently $S=A\times A\subset(\mathbb Z/m\mathbb Z)^2$
satisfies the strong progression-free condition of R-34. It contains
$(0,0),(1,0),(0,1)$, which are the three labels needed for the
irreducible specialization in that theorem. Every subset of $S$
containing those three labels retains the condition. Thus R-34
applies simultaneously as an existence statement for every
\[
                    3\leq r\leq R_m:=(|B|+2)^2.
\]
For all sufficiently large $m$, one has $L\geq m/7$ and hence
\[
 R_m\geq \frac{m^2}{49}
                \exp\!\bigl(-2C_0\sqrt{\log m}\bigr)
       \geq m^2\exp\!\bigl(-C\sqrt{\log m}\bigr)
\]
after increasing the absolute constant $C$. Also
$|A|\leq (m+1)/2<m$ for $m>1$, so $R_m<m^2$.
Finally $C\sqrt{\log m}\leq\delta\log m$ for sufficiently
large $m$, giving the stated $m^{2-\delta}$ range.
\end{proof}

\paragraph{Source and scope.}
The only additional external input is R-B19: Behrend, F.~A.
(1946). On sets of integers which contain no three terms in
arithmetical progression. \emph{Proceedings of the National Academy
of Sciences of the United States of America, 32}(12), 331--332.
\url{https://doi.org/10.1073/pnas.32.12.331}.
The original printed pages 331--332 were read in full. The new
deduction is the translated augmentation retaining $0$ and $1$,
the no-wrap strong modular condition, and its application to R-34.
The result does not give an efficient procedure for computing the
resulting representation multiplicities, and it makes no assertion
of a quadratic row bound with a fixed positive density.
